\section{Análisis}

Los conjuntos \textbf{PRIMERO} y \textbf{SIGUIENTE} son estructuras fundamentales en el diseño de
analizadores sintácticos predictivos (LL(1)). Su cálculo correcto a partir de una gramática libre
de contexto determina qué terminal puede iniciar una derivación (\textbf{PRIMERO}) y qué terminal
puede aparecer inmediatamente a la derecha de un no terminal en alguna forma sentencial
(\textbf{SIGUIENTE}).

Para el cómputo de ambos conjuntos se adoptó un enfoque iterativo de punto fijo: se aplican las
reglas de inferencia de manera repetida hasta que ningún conjunto varía entre dos iteraciones
consecutivas. Este método garantiza la convergencia para cualquier gramática libre de contexto
finita.

\subsection*{Reglas de PRIMERO}

\begin{enumerate}
  \item Si $X$ es un terminal, entonces $\text{PRIMERO}(X) = \{X\}$.
  \item Si existe la producción $X \to \varepsilon$, se añade $\varepsilon$ a $\text{PRIMERO}(X)$.
  \item Para cada producción $X \to Y_1\, Y_2\, \cdots\, Y_n$:
  \begin{itemize}
    \item Se añade $\text{PRIMERO}(Y_1) - \{\varepsilon\}$ a $\text{PRIMERO}(X)$.
    \item Si $\varepsilon \in \text{PRIMERO}(Y_i)$ para todo $i < k$, se añade también
          $\text{PRIMERO}(Y_k) - \{\varepsilon\}$.
    \item Si $\varepsilon \in \text{PRIMERO}(Y_i)$ para todo $i$, se añade $\varepsilon$ a
          $\text{PRIMERO}(X)$.
  \end{itemize}
\end{enumerate}

\subsection*{Reglas de SIGUIENTE}

\begin{enumerate}
  \item $\$ \in \text{SIGUIENTE}(S)$ donde $S$ es el símbolo inicial de la gramática.
  \item Para cada producción $A \to \alpha\, B\, \beta$:
  \begin{itemize}
    \item Se añade $\text{PRIMERO}(\beta) - \{\varepsilon\}$ a $\text{SIGUIENTE}(B)$.
    \item Si $\varepsilon \in \text{PRIMERO}(\beta)$ (o $\beta = \varepsilon$), se añade
          $\text{SIGUIENTE}(A)$ a $\text{SIGUIENTE}(B)$.
  \end{itemize}
\end{enumerate}

\vspace{1cm}

El proyecto implementa ambos algoritmos en Python mediante una arquitectura modular dividida en
tres capas: modelos, persistencia y presentación. La capa de modelos define la clase
\texttt{Gramatica}; la capa de persistencia agrupa los módulos de lectura de archivo y de cómputo
de los conjuntos; y la capa de presentación contiene la interfaz gráfica construida con Tkinter.
